\documentclass[11pt]{article}
\usepackage[margin=1in]{geometry}
\usepackage{amsmath,amssymb,amsthm}
\newtheorem{theorem}{Theorem}
\title{A terminal twelve-row frame with an optimal two-row neighbor escape}
\author{Joshua Mark}
\date{20 November 2026; finite certificate, priority unestablished}
\begin{document}\maketitle
This certificate concerns order twenty, where Hadamard matrices are
already known. It gives an exact trap for monotone row addition and
an exact minimal-loss escape in the self-dual-code construction.

\begin{theorem}[Finite certified example]
There is a self-dual ten-dimensional code $C\subseteq\mathbb F_5^{20}$
with exactly twelve antipodal classes of sign words. Representatives
with first coordinate $+1$ form a real orthogonal $12\times20$ sign
frame $P$, and span $C$ over $\mathbb F_5$.
No further sign row is orthogonal to every row of $P$.
Every new sign class is orthogonal to at most ten old rows, so no
one-row deletion admits even one genuinely new row.
Deleting the two old rows indexed $3,8$ (zero based) permits a full
Hadamard completion of the remaining ten rows. Its mod-five code
$C'$ is a self-dual neighbor: $\dim(C\cap C')=9$.
One explicit neighbor pivot takes the sign count from twelve to twenty.
Thus two is the exact minimum number of old rows lost in such an escape.
\end{theorem}

The complete basis and twelve sign representatives are stored in
\texttt{search.json}. The construction starts with the first five rows
of the Paley order-twenty matrix and extends their span over
$\mathbb F_5$ using deterministic seed $2026112000$.
The construction is only a way to find the certificate; correctness
is checked from its stored entries. The code basis is self-orthogonal
of rank ten. Enumerating its $2^9=512$ systematic sign assignments
finds exactly twelve normalized sign words, of mod-five rank ten;
their real Gram matrix is $20I_{12}$.
If a new sign row were exactly orthogonal to $P$, its residue would
belong to $C^\perp=C$, hence to the already exhausted twelve classes.
Every such class fails orthogonality to itself, proving terminality.

An independent binary enumeration checks all $2^{19}=524288$
normalized sign words using only integer Hamming distances. Excluding
the twelve existing words, its compatibility counts are
\[
\begin{array}{c|rrrrrrrrrrr}
\text{compatible}&0&1&2&3&4&5&6&7&8&9&10\\\hline
\text{new classes}&262840&6840&27092&48948&67344&58976&34920&12600&3992&616&108
\end{array}
\]
No new word is orthogonal to eleven or twelve old rows. The 108
best words group by their two nonorthogonal old rows. For the pair
$(3,8)$, all ten available words are mutually orthogonal; together
with the ten retained old words they give Gram matrix $20I_{20}$.
Their normalized negative-coordinate bit masks are
\[
7676,126672,307898,444300,581982,645928,683178,816740,855582,886242.
\]
Bit $j$ represents coordinate $j$, with bit one denoting a minus sign.
The retained old rows have modular rank nine, the completed code
rank ten, and the union of the old and new code bases rank eleven.
These prove the claimed intersection dimension. If $v$ is the first
new row, then
\[
 C'=(C\cap v^\perp)+\langle v\rangle,
\]
as independently checked by row reduction.

The scripts \texttt{search.py} and \texttt{verify.py} use only the
Python standard library and original project code from the adjacent
\texttt{codes/check.py}. The latter also writes
\texttt{verification.json} and \texttt{exchange.json}. No external
implementation is imported. This example disproves the strategy of
always retaining all previously selected rows; it does not assert
that two-row neighbor escapes suffice in general or yield unknown
Hadamard orders.

\begin{theorem}[Elementary tensor persistence]
If a sign frame $P$ admits no orthogonal sign row and $H$ is a square
Hadamard matrix of order $t$, then $P\otimes H$ also admits no
orthogonal sign row. In particular this certificate gives terminal
$12t\times20t$ frames whenever such an $H$ is supplied.
\end{theorem}
\begin{proof}
Reshape a hypothetical new sign vector into a $20\times t$ sign
matrix $X$. Orthogonality to $P\otimes H$ says $PXH^T=0$, under
the corresponding vectorization convention. Since $H$ is invertible,
$PX=0$. Each column of $X$ would therefore be a sign vector
orthogonal to $P$, a contradiction.
\end{proof}
\end{document}
